\documentclass[10pt,reqno]{amsart}

\title{Research Statement}
\author{Wenrui Huang}
\usepackage{geometry,mathtools,amssymb,amsthm,tcolorbox,esint,xcolor}
\geometry{left=1.2in,right=1.2in,top=1.1in,bottom=1.4in}
\usepackage{hyperref}

\newtheorem{theorem}{Theorem}

\theoremstyle{definition}
\newtheorem*{question}{Question}

\def\pr{\partial}
\def\H{\mathcal{H}}
\def\R{\mathbb{R}}
\def\X{\mathcal{X}}
\def\V{\mathcal{V}}

\newcommand{\rev}[1]{{\color{red}#1}}







\title{Research Statement}
\author{Wenrui Huang}
\markboth{Research Statement}{Wenrui Huang}
\begin{document}
\begin{center}
{\LARGE Research Statement}\bigskip 

{\large Wenrui Huang}
\end{center}
\vskip 22pt
 
Many physical systems involve large numbers of interacting particles or waves, for which tracking every microscopic degree of freedom is neither practical nor illuminating. In the problems I study, kinetic and dispersive partial differential equations provide effective descriptions of this collective behavior: kinetic equations describe particle distributions evolving under self-consistent forces, while dispersive equations describe the propagation and interaction of waves. A central example in my work is the Vlasov--Poisson system, which models collisionless plasmas in the repulsive case and self-gravitating systems such as galaxies in the attractive case. Once global existence of solutions is understood, two basic dynamical questions arise: how do solutions behave at large times, and are stationary solutions stable under small perturbations? My research addresses these questions through three main themes.

\begin{itemize}
   \item 
First, using canonical transformations and vector-field methods, I have obtained precise descriptions of the long-time behavior and scattering of several kinetic and dispersive PDEs, including long-range regimes in which the interaction produces persistent corrections to free motion. These results are  based on  joint work with H. Kwon and M. Xie \cite{HK24,HK26,HX26a,HX26b}  and on  my single-author work \cite{H24a,H26}.
    \item Second, combining convex geometry with elliptic analysis, I have proved the asymptotic stability of the vacuum for the Vlasov--Poisson system in a convex domain, showing how the geometry of the domain shapes the asymptotic dynamics. This is based on  joint work with B. Pausader and M. Suzuki \cite{HPS24}.
    \item Third, combining Hamiltonian dynamics with spectral and modulation analysis, I have established a criterion that determines the linear stability or instability of small BGK waves, spatially inhomogeneous stationary states of the Vlasov--Poisson system. This is based on  joint work with D. Bian, E. Grenier and B. Pausader \cite{BGHP26}.
 
\end{itemize}
The remainder of this statement develops these themes in turn and concludes with outlining ongoing projects and future directions.


\section*{Theme 1: Long-Time Behavior and Scattering for Kinetic and Dispersive Equations}
 
\noindent\textbf{Scattering theory in classical Vlasov models.} In joint work with H. Kwon and M. Xie and in my own work \cite{H26,HK24,HK26,HX26b}, I study how the interaction potential $\Gamma$ and the external potential $V$ shape the asymptotic dynamics of collisionless kinetic equations in Hamiltonian form,
\begin{equation}\label{eq:VP}
\pr_t\mu+\{\mu,\H\}=0,\qquad \H=\frac{|v|^2}{2}+V(x)-\lambda\phi,\qquad \phi=\Gamma*\rho,
\end{equation}
where $\{\cdot,\cdot\}$ is the Poisson bracket in $(x,v)$, $\mu^2$ is the phase-space density and $\rho=\int\mu^2\,dv$ is the spatial density. For $d=3$, $V\equiv0$ and $\Delta\Gamma=\delta_0$, this is the classical Vlasov--Poisson system, repulsive (plasma) for $\lambda>0$ and attractive (gravitational) for $\lambda<0$.
 
\smallskip
 
\noindent\emph{Three asymptotic regimes.} We compare the solution with free transport $(x,v)\mapsto(x+tv,v)$. For the Riesz interaction $\Gamma=|x|^{-\alpha}$ in three dimensions (Coulomb when $\alpha=1$), the density has size $t^{-3}$ on $|x|\lesssim t$, so the force is of order $t^{-1-\alpha}$. The velocity converges, but its correction, of order $t^{-\alpha}$, accumulates in the position: the deviation from free motion $\int_1^t\tau^{-\alpha}\,d\tau$ is bounded for $\alpha>1$, of order $\log t$ for $\alpha=1$, and of order $t^{1-\alpha}$ for $\alpha<1$.
\begin{itemize}
\item \emph{Linear scattering} ($\alpha>1$): $\mu(t,x+tv,v)\to\mu_\infty(x,v)$.
\item \emph{Logarithmic modified scattering} ($\alpha=1$): $\mu\bigl(t,x+tv+\log t\,E_\infty(v),v\bigr)\to\mu_\infty(x,v)$, where the asymptotic field $E_\infty$ is determined by the limiting density.
\item \emph{Long-range modified scattering} ($\alpha<1$): the correction $t^{1-\alpha}$ changes the density, and hence the force, at relative order $t^{-\alpha}$, producing further terms $t^{1-2\alpha},t^{1-3\alpha},\dots$; those with $n\alpha\leq1$ do not converge, so the trajectories carry $\lfloor1/\alpha\rfloor$ corrections, as in \eqref{eq:long-range-scatter}.
\end{itemize}
The same trichotomy holds for the harmonic potential, relative to the free hyperbolic flow, and for the Hartree equation, where the accumulated potential enters the phase. The \emph{(modified) wave operator} solves the converse problem, producing a solution with a prescribed asymptotic profile $\mu_\infty$; together with the forward result, it defines the \emph{scattering map} $\mu_{-\infty}\mapsto\mu_{+\infty}$.
 
For the classical Vlasov--Poisson system, global existence \cite{BD85,LP91,P92}, the failure of linear scattering \cite{CH11}, logarithmic modified scattering \cite{CK16,IPWW22}, and modified wave operators and scattering maps \cite{FOPW21} are known. We determine how this picture changes under a repulsive harmonic potential and under interactions decaying more slowly than Coulomb: we identify the regime and the corrected trajectories and, for the harmonic potential and for Riesz interactions near the Coulomb exponent, construct the (modified) wave operators and scattering maps.
 
\smallskip
 
\noindent\emph{Repulsive harmonic potential.} Our first theorem completes the small-data analysis of \cite{VV24,BVV25} with the construction of wave operators and scattering maps \cite{HK26} and with the case $d=1$ \cite{H26}.
\begin{theorem}[\cite{H26,HK26}]\label{thm:VPH}
Let $\lambda\in\R$, $\Delta\Gamma=\delta_0$ and $V(x)=-|x|^2/2$ in dimension $d$. For initial data sufficiently small in some weighted Sobolev space, \eqref{eq:VP} admits a unique global solution, which scatters linearly if $d\geq3$, exhibits logarithmic modified scattering if $d=2$, and exhibits long-range modified scattering if $d=1$. In each case, the (modified) wave operators exist for small data. When $d=1$ and $\lambda>0$ (repulsive), these conclusions hold for \emph{arbitrarily large} initial data.
\end{theorem}
We introduce the \emph{hyperbolic lens transform}: with $s=\tanh t$, it sends $t=\infty$ to $s=1$, absorbs the harmonic potential, and preserves the Hamiltonian structure. The transformed force carries the factor $(1-s^2)^{(d-4)/2}$, which is integrable up to $s=1$ for $d\geq3$ (linear scattering) and borderline nonintegrable for $d=2$ (logarithmic correction); in these dimensions one derivative and no spatial moments suffice, fewer than the vector-field approaches of \cite{VV24,BVV25} require. For $d=1$ the factor $(1-s^2)^{-3/2}$ is too singular, and we introduce a \emph{Hamilton--Jacobi correction}: we shift the momentum by $\pr_q\Theta$, where $\Theta$ solves a Hamilton--Jacobi equation forced by the self-consistent potential. This replaces the singular force by a transport term and, for $\lambda>0$, damps the momentum, which gives uniform bounds for data of arbitrary size.
 
\smallskip
 
\noindent\emph{Riesz interactions below the Coulomb exponent.} To our knowledge, \cite{HK24} gave the first kinetic model exhibiting modified scattering with a power-law trajectory correction; Hong and Pankavich \cite{HP26} extended forward modified scattering to $1/2<\alpha<1$. Our next theorem extends it to all $0<\alpha<1$.
\begin{theorem}[\cite{HK24,HX26b}]\label{thm:VR}
Let $\lambda\in\R$, $d=3$ and $\Gamma=|x|^{-\alpha}$ with $\alpha\in(0,1)$. For initial data sufficiently small in some weighted Sobolev space, \eqref{eq:VP} admits a unique global solution, which exhibits modified scattering: there exist bounded functions $f_\infty$ and $\theta_k$, $k=1,2,\dots,\lfloor1/\alpha\rfloor$, such that
\begin{equation}\label{eq:long-range-scatter}
\Bigl\|f(t,x,v)-f_\infty\Bigl(\frac{x}{t},\,x-tv-\sum_{n=1}^{\lfloor1/\alpha\rfloor}\frac{t^{1-n\alpha}-1}{1-n\alpha}\nabla\theta_n\Bigl(\frac{x}{t}\Bigr)\Bigr)\Bigr\|_{L^2_{x,v}\cap L^\infty_{x,v}}\lesssim t^{-\alpha}\to0.
\end{equation}
\end{theorem}
Together with the Coulomb case, Theorem~\ref{thm:VR} shows that $\alpha=1$ is the threshold between logarithmic and long-range modified scattering, and that the number $\lfloor1/\alpha\rfloor$ of corrections grows as $\alpha\to0$. Near the Coulomb exponent, \cite{HK24} uses a \emph{kinetic pseudoconformal transformation} sending $t=\infty$ to a finite endpoint, together with a symplectic change of variables that controls the endpoint singularity and yields modified wave operators and scattering maps.
 
\smallskip
 
\noindent\textbf{The Vlasov--Hartree system across the Coulomb threshold.} In joint work with M. Xie \cite{HX26a,HX26b}, we study the three-dimensional Vlasov--Hartree system, a mixed classical--quantum model of a Bose--Einstein condensate interacting with a degenerate Fermi gas, derived from many-body Bose--Fermi dynamics \cite{CMP23}:
\begin{equation}\label{eq:Vlasov-Hartree}
\begin{cases}
\partial_t f+v\cdot\nabla_x f
-\nabla_x\bigl(V_\alpha*|u|^2\bigr)\cdot\nabla_v f=0,\\[0.3em]
i\partial_tu-\Delta u=(V_\alpha*\rho_f)u,
\qquad
\displaystyle
\rho_f(t,x)=\int_{\mathbb R^3}f^2(t,x,v)\,dv,
\end{cases}
\qquad
V_\alpha(x)=\lambda|x|^{-\alpha},
\quad \lambda\neq0,
\end{equation}
where $f^2$ is the fermionic phase-space density and $u$ is the bosonic wave function. The coupling is purely between the species, so the fermionic trajectory correction and the bosonic phase correction must be constructed together.
\begin{theorem}[\cite{HX26a,HX26b}]\label{thm:VH}
Let $1/2<\alpha<2$. For initial data sufficiently small in some weighted Sobolev spaces, \eqref{eq:Vlasov-Hartree} admits a unique global strong solution, and both components exhibit long-range modified scattering in the long-range regime $0<\alpha<1$, logarithmic modified scattering at the Coulomb threshold $\alpha=1$, and linear scattering in the short-range regime $1<\alpha<2$.
\end{theorem}
In \cite{HX26a} ($1/2<\alpha<2$), the Coulomb case combines Hamiltonian characteristics for the Vlasov equation with Galilean vector fields for the Schr\"odinger equation; for $1/2<\alpha<1$, a time-dependent symplectic transformation removes the leading nonintegrable trajectory correction and a regularized effective phase makes the dispersive remainder integrable; for $\alpha>1$, the coupling is integrable along free rays.
 
\smallskip
 
\noindent\textbf{A common structure for $0<\alpha\leq1/2$.} Here forward modified scattering had remained open for the Hartree, Vlasov--Riesz and Vlasov--Hartree equations; in \cite{HX26b} we treat all three at once. Along the free rays $x=ty$ ($y$ is the velocity, or the group velocity for the Schr\"odinger flow), the kinetic and Hartree profiles evolve under the same Hamiltonian $|\eta|^2/(2t^2)+t^{-\alpha}\varphi(t,y)$, where $t^{-\alpha}\varphi$ is the potential along the ray and $\eta\mapsto-i\nabla_y$ in the Hartree case. Hence a single nonlinear action $\theta(t,y)$ corrects both components, by a momentum translation $\nabla_y\theta$ on the kinetic profile and by a phase modulation on the wave. Cancelling the potential alone does not suffice: the momentum shift itself produces a quadratic term of order $t^{-2\alpha}$, nonintegrable for $\alpha\leq1/2$, so $\theta$ must solve a Hamilton--Jacobi equation. The action in turn changes the density, and hence the potential, at the same order $t^{-2\alpha}$, so the correction and the density feed back on each other. With the rescaled densities $m_{\mathrm K}(t,y)=t^3\rho_f(t,ty)$ and $m_{\mathrm H}(t,y)=t^3|u(t,ty)|^2$, each component $k\in\{\mathrm K,\mathrm H\}$ satisfies the \emph{action--density system}
\[
\pr_t\theta_k-\frac{|\nabla_y\theta_k|^2}{2t^2}=t^{-\alpha}\varphi_j,\qquad
\pr_t m_k-\frac{1}{t^2}\nabla_y\cdot\bigl(m_k\nabla_y\theta_k\bigr)=-\frac{1}{t^2}\nabla_y\cdot\mathcal J_k,
\]
where $t^{-\alpha}\varphi_j$ is the (regularized) potential generated by the density $m_j$ of the source component, with $j=k$ for the Hartree and Vlasov--Riesz equations and $j\neq k$ for the Vlasov--Hartree system, and $\mathcal J_k$ is an integrable residual current. The density drives the action, whose gradient carries the leading density flux; iterating this system produces the corrections $t^{1-\alpha},t^{1-2\alpha},\dots$ in Theorem~\ref{thm:VR}, each determined recursively from the limiting density, and the same recursion applies to the coupled system. Thus trajectory corrections in kinetic theory and phase corrections in dispersive equations are two realizations of one structure.








 
 


\section*{Theme 2: Asymptotic stability of the vacuum for the Vlasov--Poisson system in a convex domain}

Plasma--wall interactions are central to plasma physics and arise in applications ranging from neon lighting and semiconductor fabrication to fusion devices. Mathematically, the vacuum state, in which no particles are present, is the simplest equilibrium. Its asymptotic stability asks whether a sufficiently dilute plasma remains globally regular and approaches a simpler transport regime at large times. In joint work with Benoit Pausader and Masahiro Suzuki \cite{HPS24}, we studied this question for the three-dimensional Vlasov--Poisson system in a class of unbounded convex domains.

The wall is assumed to be perfectly conducting, so the electrostatic potential satisfies a Dirichlet boundary condition, while particles are absorbed when they reach the wall. Writing the phase-space particle density as $\widetilde{\mu}^{\,2}$, the system takes the form
\begin{equation}\label{VPW}
\begin{cases}
\partial_t\widetilde{\mu}
+v\cdot\nabla_x\widetilde{\mu}
-\lambda\nabla_x\phi\cdot\nabla_v\widetilde{\mu}=0,
& (x,v)\in D\times\mathbb{R}^3,\\[0.2em]
\displaystyle
\Delta_x\phi=\int_{\mathbb{R}^3}\widetilde{\mu}^{\,2}(x,v,t)\,dv,
& x\in D,\\[0.4em]
\widetilde{\mu}=0 \quad\text{on }\Gamma_{\mathrm{in}},
\qquad
\phi=0 \quad\text{on }\partial D,
\qquad
\phi(x,t)\to0 \quad\text{as }|x|\to\infty .
\end{cases}
\end{equation}
Here
\[
\Gamma_{\mathrm{in}}
=\{(x,v):x\in\partial D,\ {\bf n}(x)\cdot v\geq0\},
\]
where ${\bf n}$ denotes the inward unit normal. Thus the condition on $\Gamma_{\mathrm{in}}$ prevents particles from entering the domain through the wall and represents one-way absorption. We focus on the repulsive plasma case $\lambda<0$.

The presence of the wall creates two difficulties absent in the whole-space problem: translation invariance is lost, and the transport boundary is characteristic along grazing trajectories. Earlier work established global existence for kinetic equations in various domains \cite{G94,G95,HV10,CI23,F18}, but did not determine the long-time behavior of the surviving particles or the self-consistent field.

A central geometric object in our analysis is the asymptotic domain. After translating coordinates so that $0\in\partial D$, we define
\[
D_\infty
:=\operatorname{Int}\bigcap_{s>0}s^{-1}D.
\]
Its closure is the maximal solid convex cone with vertex at the origin contained in $\overline{D}$. Geometrically, it encodes the velocity directions that remain asymptotically compatible with escape to infinity.

\begin{theorem}\label{thm:VP-Wall}
Let $D\subset\mathbb{R}^3$ be an admissible $C^3$ unbounded convex domain. For sufficiently small, regular, and localized initial data satisfying the natural boundary compatibility conditions, the system \eqref{VPW} admits a unique global classical solution. Moreover, the electrostatic potential and electric field decay at the optimal rates $t^{-1}$ and $t^{-2}$, respectively.

If $D_\infty\neq\varnothing$, then there exist a limiting distribution $\mu_\infty$ and an asymptotic field $E_\infty$ such that
\[
\widetilde{\mu}
\bigl(x+tv+\lambda(\log t)E_\infty(v),v,t\bigr)
\longrightarrow \mu_\infty(x,v)
\]
 for $(x,v)\in\mathbb{R}^3\times D_\infty$. Here $E_\infty$ is determined jointly by $\mu_\infty$ and the Dirichlet Green function of $D_\infty$. In addition, the phase-space mass carried by velocities outside $\overline{D_\infty}$ tends to zero.
\end{theorem}

This result shows that the geometry of the wall affects not only which particles survive, but also the leading asymptotic field and the logarithmic correction to free transport. In the extreme case $D_\infty=\varnothing$, the total mass of the surviving electrons tends to zero as $t\to\infty$.

The proof combines convex geometry, elliptic theory, and kinetic transport. We first extend the Dirichlet potential across the wall, allowing the boundary problem to be analyzed through a whole-space evolution while retaining the one-way absorption mechanism. In free-transport coordinates, weighted phase-space estimates are combined with a scale decomposition of the Dirichlet Green function. The blow-down
\(
s^{-1}D\rightarrow D_\infty
\)
then identifies the leading field along an escaping trajectory as
\[
E(tv,t)\sim t^{-2}E_\infty(v).
\]
Integrating this borderline force produces the logarithmic correction in the particle trajectories. We also introduce a grazing function that quantifies proximity to tangential impact and provides uniform control near the characteristic part of the boundary. Together, the asymptotic domain and the grazing function form a framework for studying more general plasma--wall interactions.


\section*{Theme 3: Stability and Instability of Small BGK Waves}
\providecommand{\HM}{\mathrm{HM}}

Bernstein--Greene--Kruskal (BGK) waves, introduced in \cite{BGK57}, are spatially inhomogeneous steady states of the Vlasov--Poisson system. They play a central role in its long-time dynamics. Numerical simulations suggest that finite-amplitude perturbations of a stable homogeneous equilibrium, such as a Maxwellian, often evolve towards a BGK wave or a superposition of BGK waves, rather than converging back to a homogeneous state.

We consider the one-dimensional Vlasov--Poisson system
\begin{equation}\label{VP}
\partial_t\mu+\{\mu,\mathcal H\}=0,
\qquad
\mathcal H=\frac{v^2}{2}-\phi(x,t),
\qquad
\Delta\phi=\int_{\mathbb R}\mu\,dv-1,
\end{equation}
where $\{f,g\}=\partial_xf\,\partial_vg-\partial_vf\,\partial_xg$. Here $\mu(x,v,t)$ is the distribution function of electrons moving against a fixed background of ions with unit density, while $E(x,t)=-\partial_x\phi(x,t)$ is the self-consistent electrostatic field. We consider solutions that are periodic in $x$ with period $2\pi$.

For the class of steady states considered here, the particle distribution depends only on the particle energy:
\[
\mu_S(x,v)=F(\mathcal H_S(x,v)),
\qquad
\mathcal H_S(x,v)=\frac{v^2}{2}-\phi_S(x).
\]
When $\phi_S\equiv0$, this gives a homogeneous equilibrium $\mu_h(v)=F(v^2/2)$. When $\phi_S\not\equiv0$, it gives a BGK wave. After normalizing so that $\min\phi_S=0$, the level set $\{\mathcal H_S=0\}$ is the separatrix between trapped particles (negative energy) and free particles (positive energy), and we call $F|_{(0,\infty)}$ the free part of the energy distribution.

The stability theory for homogeneous equilibria is comparatively well developed. The linearization at $\mu_h$ decouples into spatial Fourier modes, and the mode $e^{i\xi x}$ is governed by the dielectric function
\begin{equation}\label{HM}
\HM[\mu_h](\theta,\xi)=1-\frac1{\xi^2}\int_{\mathbb R}\frac{\mu_h'(v)}{v+\theta/\xi}\,dv,
\qquad \xi\in\mathbb Z\setminus\{0\},\quad \operatorname{Im}\theta<0,
\end{equation}
extended to real $\theta$ by the Plemelj formula. The mode carries a growing solution exactly when $\HM[\mu_h](\cdot,\xi)$ vanishes in the open lower half-plane $\mathbb C_-$ ( Landau \cite{L46}, Penrose \cite{P60}). When it does not vanish in $\overline{{\mathbb{C}_-}}$, it is called Penrose-stable. For Penrose-stable equilibrium, Mouhot and Villani \cite{MV11} first established nonlinear Landau damping on the torus, with subsequent simplifications and refinements in \cite{BMM16,GNR21,IPWW24b}. By contrast, the stability of spatially inhomogeneous BGK waves remains largely open, even in the small-amplitude regime. Previous results \cite{D19,GL17,HM24,W26} concern special energy distributions or inhomogeneous ion distributions with well-prepared perturbations.

In joint work with D. Bian, E. Grenier, and B. Pausader \cite{BGHP26}, we study small BGK waves, that is, BGK waves in a neighborhood of the homogeneous equilibria, and we show that their stability is decided by the nearest homogeneous equilibrium alone. If BGK waves $(\mu_n,\phi_n)$ have amplitudes $A_n:=\max\phi_n\to0$ and convergent free parts, then the limiting homogeneous equilibrium $\mu_0(v)=F_0(v^2/2)$ is \emph{marginal}, $\HM[\mu_0](0,\xi_0)=0$ for some $\xi_0\geq1$, and the potentials have a universal shape: after a translation,
\[
\phi_n(x)=\frac{A_n}{2}\bigl[1+\cos(\xi_0x)\bigr]+O(A_n^2).
\]
We call $\mu_0$ a \emph{generic bifurcation germ} if, in addition, it is nondegenerate, $\mu_0''(0)\neq0$, and has no other neutral mode, that is, $\HM[\mu_0](\theta,\xi)\neq0$ for every real $\theta\neq0$ and every $\xi\neq0$. We assume $\xi_0=1$, so that the bifurcating waves have one well per period. The small BGK waves bifurcating from $\mu_0$ are those in the box $\mathcal B_\varepsilon(\mu_0)$: the BGK waves of amplitude $A\leq\varepsilon$ whose free part is $\varepsilon$-close to $\mu_0$, with uniform bounds across the separatrix.

{Crucially, even after fixing the period and the position of the potential well, the modulation space is {\emph {infinite-dimensional}}: varying the energy distribution yields a local manifold of stationary states with infinitely many tangent directions. This differs sharply from soliton modulation in dispersive equations, which typically tracks only finitely many parameters. Thus, stable perturbations require separating this stationary component and the two symmetry directions from the genuinely dynamical part. This is the sense in which stability holds \emph{up to infinite-dimensional modulation}.}

For a fixed BGK wave $(\mu_\varepsilon,\phi_\varepsilon)$, writing $\mu=\mu_\varepsilon+f$ gives the linearized system
\begin{equation}\label{LVP}
\partial_tf+\{f,\mathcal H_\varepsilon\}-\{\mu_\varepsilon,\psi\}=0,
\qquad
\Delta\psi=\int_{\mathbb R}f\,dv.
\end{equation}

\begin{theorem}\label{thm:BGK}
Let $\mu_0$ be a generic bifurcation germ with $\xi_0=1$, and consider the spectral condition
\[
\textup{(SC)}\qquad \exists\,\xi\in\mathbb Z\setminus\{0\},\ \exists\,\theta_*\in\mathbb C_-:\quad \HM[\mu_0](\theta_*,\xi)=0 .
\]
There is $\varepsilon_0=\varepsilon_0(\mu_0)>0$ such that every BGK wave $(\mu_\varepsilon,\phi_\varepsilon)\in\mathcal B_\varepsilon(\mu_0)$ with $\varepsilon\leq\varepsilon_0$ satisfies the following alternative.
\begin{itemize}
\item \emph{(\bf{Instability}.)} If \textup{(SC)} holds, then $(\mu_\varepsilon,\phi_\varepsilon)$ is spectrally and linearly unstable. More precisely, there exist $\lambda_\varepsilon\in\mathbb C$ and nonzero functions $\widehat f_\varepsilon$ and $\widehat\psi_\varepsilon$ such that
\[
\bigl(f(x,v,t),\psi(x,t)\bigr)=e^{\lambda_\varepsilon t}\bigl(\widehat f_\varepsilon(x,v),\widehat\psi_\varepsilon(x)\bigr)
\]
solves \eqref{LVP}, with
\[
\operatorname{Re}\lambda_\varepsilon=\max\bigl\{|\operatorname{Im}\theta_*|:\ \HM[\mu_0](\theta_*,\xi)=0,\ \xi\neq0,\ \theta_*\in\mathbb C_-\bigr\}+o(1)\qquad(\varepsilon\to0).
\]

\item \emph{\bf{(Stability up to modulation.})} If \textup{(SC)} fails, then $(\mu_\varepsilon,\phi_\varepsilon)$ is linearly stable up to modulation. More precisely, every solution of \eqref{LVP} whose initial data satisfy
\[
\langle v\rangle^2f_{\mathrm{in}}\in L^\infty(\mathbb T\times\mathbb R),\qquad \iint f_{\mathrm{in}}(x,v)\,dxdv=0
\]
admits a decomposition
\[
\begin{aligned}
f(x,v,t)&=f^{\mathrm{bgk}}(x,v)+f^{\mathrm{gal}}(x,v,t)+f^{\mathrm{odd}}(x,v)+f^{\mathrm{dyn}}(x,v,t),\\
\psi(x,t)&=\psi^{\mathrm{bgk}}(x)+\psi^{\mathrm{gal}}(x,t)+\psi^{\mathrm{dyn}}(x,t).
\end{aligned}
\]
Here the components have the following interpretations:
\begin{itemize}
\item $(f^{\mathrm{bgk}},\psi^{\mathrm{bgk}})$ is a time-independent weak solution of \eqref{LVP}, tangent to the manifold of BGK waves at $(\mu_\varepsilon,\phi_\varepsilon)$. It represents an infinitesimal deformation of $(\mu_\varepsilon,\phi_\varepsilon)$ through nearby stationary BGK waves;
\item $(f^{\mathrm{gal}},\psi^{\mathrm{gal}})$ is a boosted steady solution, belonging to the two-dimensional space generated by spatial translations and Galilean boosts: for some $(a,b)\in\mathbb R^2$,
\[
f^{\mathrm{gal}}(x,v,t)=\bigl[a\partial_x+b(\partial_v+t\partial_x)\bigr]\mu_\varepsilon(x,v),
\qquad
\psi^{\mathrm{gal}}(x,t)=(a+bt)\,\partial_x\phi_\varepsilon(x);
\]
\item $f^{\mathrm{odd}}=\mathrm{sgn}(v)\,g(\mathcal H_\varepsilon)\mathbf 1_{\mathcal H_\varepsilon>0}$ is a time-independent solution with identically vanishing density and electric field, which describes an imbalance between right- and left-moving free particles; it vanishes when $f_{\mathrm{in}}$ is even in $v$;
\item $(f^{\mathrm{dyn}},\psi^{\mathrm{dyn}})$ is the genuinely time-dependent component, and its electric field $E^{\mathrm{dyn}}=-\partial_x\psi^{\mathrm{dyn}}$ satisfies
\[
\int_0^\infty\bigl\|E^{\mathrm{dyn}}(\cdot,t)\bigr\|_{L^2_x}^2\,dt<\infty.
\]
\end{itemize}
\end{itemize}
\end{theorem}

{The theorem reads the stability of a small BGK wave on its germ: the wave is unstable exactly when the germ is Penrose unstable, with the growth rate of the Penrose root, and stable up to modulation otherwise.

There are two main analytical difficulties. The first comes from the qualitative change in particle dynamics across the separatrix. For homogeneous equilibria, spatial Fourier modes decouple and stability reduces to the scalar dispersion relation \eqref{HM}. For BGK waves, spatial inhomogeneity and the coexistence of trapped and passing particles couple these modes. We introduce action--angle coordinates adapted to the particle trajectories and replace the scalar dispersion relation by an operator-valued one, $\mathbf M(\theta)$ on $L^2(\mathbb T)$, whose kernel for $\operatorname{Im}\theta<0$ gives the growing modes. Since all small BGK waves have, to first order, the Hamiltonian of the nonlinear pendulum, this operator admits a universal expansion
\[
\mathbf M(\theta)=\mathbf M^{\mathrm{hom}}(\theta)+\partial_eF(0^+)\sqrt \varepsilon\;\mathbf M^{\mathrm{pend}}\bigl(\theta/\sqrt \varepsilon\bigr)+O_{L^2\to L^2}(\varepsilon),
\]
in which $\mathbf M^{\mathrm{hom}}$ is the dispersion operator of the homogeneous profile $F(v^2/2)$ built from the free part of the wave, diagonal in the Fourier basis with entries \eqref{HM}, and $\mathbf M^{\mathrm{pend}}$ depends only on the trajectories of the nonlinear pendulum. The expansion exhibits a sharp transition at the trapping scale $|\theta|\simeq\sqrt \varepsilon$, where the trapped particles can compete with the free ones.

The second difficulty is a degeneracy analogous to that encountered in the stability theory of ground states and solitary waves. BGK waves are not isolated equilibria but belong to a manifold of nearby stationary states, and the bifurcation condition forces the zero-frequency dispersion operator $\mathbf M(0)$ to have a nontrivial kernel associated with directions tangent to this manifold: the translation mode is an exact zero and the amplitude direction carries an eigenvalue of order $\sqrt \varepsilon$. Building on our analysis of the local structure of the BGK manifold, we develop a modulation analysis that separates these neutral stationary directions from the genuinely dynamical part of the perturbation. On the two-dimensional block they span, a Schur reduction expresses the dispersion relation at the trapping scale through two universal functions of the nonlinear pendulum, which do not vanish in the closed lower half-plane for either sign of $\partial_eF(0^+)$. Hence the trapped particles never produce a growing mode, and the spectral question is decided at $|\theta|\simeq 1$, by the germ: a zero of $\HM[\mu_0]$ in $\mathbb C_-$ survives as singular points of $\mathbf M$, while its absence yields the damping of the dynamical electric field.

The evenness assumption can be removed. A general one-well steady state has the form $\mu=F(\mathcal H_S)+\mathrm{sgn}(v)\,G(\mathcal H_S)\mathbf 1_{\mathcal H_S>0}$, whose odd part $G$ carries no density but a current: the stationary theory is unchanged, and $G$ contributes exactly one new, explicit block to $\mathbf M(\theta)$, which modifies the two-dimensional block above by a term of the form $-\widetilde c^{\,2}\vartheta^2$, $\vartheta=\theta/\sqrt \varepsilon$, $\widetilde c=\int_0^\infty G'(e)e^{-1}de$. The alternative of Theorem~\ref{thm:BGK} persists for the dielectric function of the full profile, whose zeros are no longer symmetric under $\theta\mapsto-\overline\theta$: the growing modes in general travel, and in the stable case the centering of the wave is displaced by the odd part of the perturbation, by an amount proportional to $\widetilde c$.





\section*{Ongoing Projects and Future Research}

The results described above suggest several natural directions for further research.

\subsection*{Asymptotic Dynamics of the Vlasov--Poisson System with External Potentials}

\begin{enumerate}

\item \textit{Attractive harmonic potentials.}
 Theorem \ref{thm:VPH} concerns a repulsive harmonic potential. The corresponding dynamics for an attractive harmonic potential remain largely unexplored. Our lens transform also applies in the attractive setting and may provide a useful starting point. At present, only a few results are available for attractive external potentials \cite{CL22,CL24,KW25}. 

\item \textit{More general external potentials.}
A broader goal is to identify the radial external potentials for which solutions exhibit scattering or modified scattering. For radially symmetric solutions, the Hamiltonian dynamics reduce to an explicitly integrable system, and the asymptotic action--angle variables developed in \cite{PW21,PWY22} may be useful. The problem is substantially harder for nonradial solutions, even when the external potential remains radial, because this reduction is no longer available. In this setting, the Coulomb and harmonic potentials are distinguished as the two classical central potentials with superintegrable Hamiltonian flows, a structure that plays an essential role in \cite{PWY22}.

\end{enumerate}

\subsection*{Long-Time Behavior and Asymptotic Stability for Vlasov Models}

\begin{enumerate}

\item \textit{Reflecting boundaries.}
Our work \cite{HPS24} considers an absorbing boundary condition, the simplest model of plasma--wall interaction. A natural next step is to study a perfectly reflecting boundary, for which particles undergo specular reflection rather than being absorbed. The goal is to understand how repeated interactions with the boundary affect global regularity, asymptotic stability, and the long-time distribution of particles.

\end{enumerate}

\subsection*{Stability of Small BGK Waves}

\begin{enumerate}

\item \textit{Quantitative damping and nonlinear stability.}
For the linearly stable BGK waves in Theorem \ref{thm:BGK}, our current result yields a qualitative form of Landau damping: the dynamical electric field satisfies
\(
E^{\mathrm{dyn}}
\in
L^2\bigl([0,\infty);L_x^2\bigr).
\)
We are working toward quantitative damping estimates and, under suitable regularity assumptions, exponential decay of the electric field. A longer-term goal is to establish nonlinear stability.

\item \textit{Nonlinear instability.}
In the unstable regime of Theorem \ref{thm:BGK}, the next step is  to promote the exponentially growing linear mode to a nonlinear instability result, using ideas developed in \cite{GS95,L01,L05}.

\item \textit{The degenerate threshold.}
The case
\(
\partial_e F(0,0)=0
\)
is not covered by the sign criterion in Theorem \ref{thm:BGK}. Understanding this degenerate regime may require a higher-order expansion of the Hamiltonian $\mathcal H_\varepsilon$ and of the associated dispersion operator.

\end{enumerate}






\begingroup
\setlength{\parskip}{0pt}
\begin{thebibliography}{CSFGKA13}
\fontsize{9pt}{10.3pt}\selectfont
\setlength{\itemsep}{1pt}
\setlength{\parsep}{0pt}




\bibitem[BD85]{BD85}
C. Bardos and P. Degond,
\emph{Global existence for the Vlasov--Poisson equation in $3$ space variables with small initial data},
Ann. Inst. H. Poincar\'e Anal. Non Lin\'eaire \textbf{2} (1985), no.~2, 101--118.

\bibitem[BBDR05]{BBDR05}
J.~Barr\'e, F.~Bouchet, T.~Dauxois, and S.~Ruffo,
\emph{Large deviation techniques applied to systems with long-range interactions},
J. Stat. Phys. \textbf{119} (2005), nos.~3--4, 677--713.



\bibitem[BMM16]{BMM16}
J. Bedrossian, N. Masmoudi, and C. Mouhot,
\href{https://link.springer.com/article/10.1007/s40818-016-0008-2}{\emph{Landau damping: Paraproducts and Gevrey regularity}},
Ann. PDE \textbf{2} (2016), no.~1, Art.~4, 71 pp.

\bibitem[BGK57]{BGK57}
I. Bernstein, J. Greene, and M. Kruskal,
\href{https://journals.aps.org/pr/abstract/10.1103/PhysRev.108.546}{\emph{Exact nonlinear plasma oscillations}},
Phys. Rev. \textbf{108} (1957), no.~3, 546--550.

\bibitem[BGHP26]{BGHP26}
D. Bian, E. Grenier, W. Huang, and B. Pausader,
\href{https://arxiv.org/abs/2601.10030}{\emph{Stability and instability of small BGK waves}},
arXiv:2601.10030.




\bibitem[BVV25]{BVV25}
L.~Bigorgne, A.~Velozo Ruiz, and R.~Velozo Ruiz,
\emph{Modified scattering of small data solutions for the Vlasov--Poisson system with a repulsive potential},
SIAM J. Math. Anal. \textbf{57} (2025), no.~1, 714--752.







%\bibitem[CMMP25]{CMMP25}
%E. C\'ardenas, J.~K. Miller, D. Mitrouskas, and N. Pavlovi\'c,
%\href{https://arxiv.org/abs/2502.18678}{\emph{Emergence of fermion-mediated interactions in Bose--Fermi mixtures}},
%arXiv:2502.18678 (2025).

\bibitem[CMP23]{CMP23}
E. C\'ardenas, J.~K. Miller, and N. Pavlovi\'c,
\href{https://arxiv.org/abs/2309.04638}{\emph{On the effective dynamics of Bose--Fermi mixtures}},
Ars Inven. Anal. (2025), Paper No.~5, 54 pp.; \href{https://arxiv.org/abs/2309.04638}{arXiv:2309.04638} (2023).




\bibitem[CI23]{CI23}
L. Cesbron and M. Iacobelli,
\href{https://doi.org/10.2140/apde.2023.16.2465}{\emph{Global well-posedness of Vlasov--Poisson-type systems in bounded domains}},
Anal. PDE \textbf{16} (2023), no.~10, 2465--2494.




\bibitem[CL22]{CL22}
S. Chaturvedi and J. Luk,
\href{https://www.aimsciences.org/article/doi/10.3934/krm.2022002}{\emph{Phase mixing for solutions to 1D transport equation in a confining potential}},
Kinet. Relat. Models \textbf{15} (2022), no.~3, 403--416.

\bibitem[CL24]{CL24}
S.~Chaturvedi and J.~Luk,
\emph{Linear and nonlinear phase mixing for the gravitational Vlasov--Poisson system under an external Kepler potential},
to appear in Arch. Ration. Mech. Anal (2026).

\bibitem[CH11]{CH11}
S.-H.~Choi and S.-Y.~Ha,
\emph{Asymptotic behavior of the nonlinear Vlasov equation with a self-consistent force},
SIAM J. Math. Anal. \textbf{43} (2011), no.~5, 2050--2077.



\bibitem[CK16]{CK16}
S.-H.~Choi and S.~Kwon,
\emph{Modified scattering for the Vlasov--Poisson system},
Nonlinearity \textbf{29} (2016), no.~9, 2755--2774.





\bibitem[D19]{D19}
B. Despr\'es,
\href{https://hal.sorbonne-universite.fr/hal-01985113v1/document}{\emph{Scattering structure and Landau damping for linearized Vlasov equations with inhomogeneous Boltzmannian states}},
Ann. Henri Poincar\'e \textbf{20} (2019), 2767--2818.



\bibitem[F18]{F18}
X. Fern\'andez-Real,
\emph{The Lagrangian structure of the Vlasov--Poisson system in domains with specular reflection},
Comm. Math. Phys. \textbf{364} (2018), 1327--1406.

\bibitem[FOPW21]{FOPW21}
P. Flynn, Z. Ouyang, B. Pausader, and K. Widmayer,
\href{https://doi.org/10.1007/s42543-021-00041-x}{\emph{Scattering map for the Vlasov--Poisson system}},
Peking Math. J. \textbf{6} (2023), no.~2, 365--392; \href{https://arxiv.org/abs/2101.01390}{arXiv:2101.01390} (2021).





\bibitem[GNR21]{GNR21}
E. Grenier, T.~T. Nguyen, and I. Rodnianski,
\emph{Landau damping for analytic and Gevrey data},
Math. Res. Lett. \textbf{28} (2021), no.~6, 1679--1702.


\bibitem[G94]{G94}
Y. Guo,
\emph{Regularity for the Vlasov equations in a half-space},
Indiana Univ. Math. J. \textbf{43} (1994), 255--320.

\bibitem[G95]{G95}
Y. Guo,
\emph{Singular solutions for the Vlasov--Maxwell system on a half line},
Arch. Ration. Mech. Anal. \textbf{131} (1995), 241--304.

\bibitem[GL17]{GL17}
Y. Guo and Z. Lin,
\href{https://link.springer.com/article/10.1007/s00220-017-2873-2}{\emph{The existence of stable BGK waves}},
Comm. Math. Phys. \textbf{352} (2017), 1121--1152.

\bibitem[GS95]{GS95}
Y. Guo and W.~A. Strauss,
\href{https://onlinelibrary.wiley.com/doi/abs/10.1002/cpa.3160480803}{\emph{Instability of BGK equilibria}},
Comm. Pure Appl. Math. \textbf{48} (1995), no.~8, 861--894.

\bibitem[HM24]{HM24}
M. Had\v{z}i\'c and M. Moreno,
\href{https://arxiv.org/abs/2412.07025}{\emph{On absence of embedded eigenvalues and stability of BGK waves}},
Ann. Henri Poincar\'e (2026).


\bibitem[HP26]{HP26}
Y. Hong and S. Pankavich,
\href{https://arxiv.org/abs/2604.04256}{\emph{Modified scattering for the Vlasov--Riesz system with long-range interactions}},
arXiv:2604.04256 (2026).




\bibitem[H24a]{H24a}
W. Huang,
\href{https://arxiv.org/abs/2411.02737}{\emph{Modified wave operators for the Hartree equation with repulsive Coulomb potential}},
arXiv:2411.02737 (2024).

\bibitem[H26]{H26}
W. Huang
\href{https://arxiv.org/abs/2610.02551}{\emph{Asymptotic behavior of the one-dimensional Vlasov--Poisson system with a repulsive harmonic potential}},



\bibitem[HK24]{HK24}
W. Huang and H. Kwon,
\href{https://link.springer.com/article/10.1007/s00023-026-01735-7}{\emph{Scattering of the Vlasov--Riesz system in three dimensions}},
Ann. Henri Poincar\'e (2026), \href{https://arxiv.org/abs/2407.16919}{arXiv:2407.16919}.

\bibitem[HK26]{HK26}
W. Huang and H. Kwon,
\href{https://arxiv.org/abs/2602.21344}{\emph{Scattering map for the Vlasov--Poisson system with a repulsive harmonic potential}},
arXiv:2602.21344 (2026).

\bibitem[HPS24]{HPS24}
W. Huang, B. Pausader, and M. Suzuki,
\href{https://arxiv.org/abs/2412.13434}{\emph{The Vlasov--Poisson system with a perfectly conducting wall: Convex domains}},
 Comm. Math. Phys. \textbf{407} (2026)

\bibitem[HX26a]{HX26a}
W. Huang and M. Xie,
\href{https://arxiv.org/abs/2608.04402}{\emph{Long-time dynamics of Vlasov--Hartree systems across the Coulomb threshold}},
arXiv:2608.04402 (2026).

\bibitem[HX26b]{HX26b}
W. Huang and M. Xie,
\href{https://arxiv.org/abs/2610.02643}{A common structure for modified scattering: Vlasov--Riesz, Hartree, and their coupling},


\bibitem[HV10]{HV10}
H.~J. Hwang and J.~J. Vel\'azquez,
\emph{Global existence for the Vlasov--Poisson system in bounded domains},
Arch. Ration. Mech. Anal. \textbf{195} (2010), 763--796.

\bibitem[IRW24]{IRW24}
M.~Iacobelli, S.~Rossi, and K.~Widmayer,
\emph{On the stability of vacuum in the screened Vlasov--Poisson equation},
J. Lond. Math. Soc. \textbf{113} (2026), no.~1, e70426; \href{https://arxiv.org/abs/2410.17978}{arXiv:2410.17978} (2024).



\bibitem[IPWW22]{IPWW22}
A.~D. Ionescu, B. Pausader, X. Wang, and K. Widmayer,
\href{https://doi.org/10.1093/imrn/rnab155}{\emph{On the asymptotic behavior of solutions to the Vlasov--Poisson system}},
Int. Math. Res. Not. IMRN \textbf{2022} (2022), no.~12, 8865--8889.





\bibitem[IPWW24b]{IPWW24b}
A.~D. Ionescu, B. Pausader, X. Wang, and K. Widmayer,
\emph{Nonlinear Landau damping and wave operators in sharp Gevrey spaces},
preprint (2024), \href{https://arxiv.org/abs/2405.04473}{arXiv:2405.04473}.











\bibitem[KW25]{KW25}
B. Kepka and K. Widmayer,
\href{https://arxiv.org/abs/2511.04363}{\emph{Modified scattering dynamics in the Vlasov--Poisson equation near an attractive point mass}},
arXiv:2511.04363 (2025).

\bibitem[L46]{L46}
L. Landau,
\emph{On the vibrations of the electronic plasma},
Akad. Nauk SSSR. Zhurnal Eksper. Teoret. Fiz. \textbf{16} (1946), 574--586; English transl.\ J. Phys. (USSR) \textbf{10} (1946), 25--34.




\bibitem[LMR05]{LMR05}
M. Lemou, F. M\'ehats, and P. Rapha\"el,
\emph{Orbital stability and singularity formation for Vlasov--Poisson systems},
C. R. Acad. Sci. Paris, Ser. I \textbf{341} (2005), 269--274.










\bibitem[L01]{L01}
Z. Lin,
\href{https://link.intlpress.com/JDetail/1806606498308452354}{\emph{Instability of periodic BGK waves}},
Math. Res. Lett. \textbf{8} (2001), no.~4, 521--534.

\bibitem[L05]{L05}
Z. Lin,
\href{https://onlinelibrary.wiley.com/doi/abs/10.1002/cpa.20028}{\emph{Nonlinear instability of periodic BGK waves for the Vlasov--Poisson system}},
Comm. Pure Appl. Math. \textbf{58} (2005), no.~4, 505--528.





\bibitem[LP91]{LP91}
P.-L.~Lions and B.~Perthame,
\emph{Propagation of moments and regularity for the $3$-dimensional Vlasov--Poisson system},
Invent. Math. \textbf{105} (1991), no.~2, 415--430.






\bibitem[MV11]{MV11}
C. Mouhot and C. Villani,
\href{https://projecteuclid.org/journals/acta-mathematica/volume-207/issue-1/On-Landau-damping/10.1007/s11511-011-0068-9.full}{\emph{On Landau damping}},
Acta Math. \textbf{207} (2011), 29--201.





\bibitem[PW21]{PW21}
B. Pausader and K. Widmayer,
\href{https://link.springer.com/article/10.1007/s00220-021-04117-8}{\emph{Stability of a point charge for the Vlasov--Poisson system: The radial case}},
Comm. Math. Phys. \textbf{385} (2021), no.~3, 1741--1769.

\bibitem[PWY22]{PWY22}
B. Pausader, K. Widmayer, and J. Yang,
\href{https://doi.org/10.4171/JEMS/1518}{\emph{Stability of a point charge for the repulsive Vlasov--Poisson system}},
J. Eur. Math. Soc. \textbf{28} (2026), 2751--2848.

\bibitem[P60]{P60}
O. Penrose,
\emph{Electrostatic instabilities of a uniform non-Maxwellian plasma},
Phys. Fluids \textbf{3} (1960), 258--265.




\bibitem[P92]{P92}
K.~Pfaffelmoser,
\emph{Global classical solutions of the Vlasov--Poisson system in three dimensions for general initial data},
J. Differential Equations \textbf{95} (1992), no.~2, 281--303.





\bibitem[S20]{S20}
S.~Serfaty,
\emph{Mean field limit for Coulomb-type flows},
Duke Math. J. \textbf{169} (2020), no.~15, 2887--2935, with an appendix by M.~Duerinckx and S.~Serfaty.





\bibitem[VV24]{VV24}
A.~Velozo Ruiz and R.~Velozo Ruiz,
\emph{Small data solutions for the Vlasov--Poisson system with a repulsive potential},
Comm. Math. Phys. \textbf{405} (2024), no.~3, Paper No.~80, 45 pp.

\bibitem[W26]{W26}
R.~Weder,
\newblock \emph{Asymptotic stability of Bernstein--Greene--Kruskal (BGK) waves, Landau damping, and scattering theory},
\newblock preprint, arXiv:2606.16194, 2026.


\bibitem[W24]{W24}
D.~Wei,
\href{https://doi.org/10.3934/cpaa.2025054}{\emph{Nonlinear stability of the one dimensional screened Vlasov--Poisson equation}},
Comm. Pure Appl. Anal. \textbf{31} (2026), 65--85; 

\end{thebibliography}
\endgroup

\end{document}